\par\noindent\begin{minipage}{\linewidth}
% !TeX root = ../教材题库_全部条目_转写工作稿.tex
% Source: XJ-B1-C02-TH-P050-Y1033
% 第2章 / 2.3 一元二次不等式 / 教材 p.50 / PDF 第57页
% 已对照原页转写题干及公式；待教师复核；未制作答案
% 原题及相邻原解截图（核对用）：../题目原图/XJ-B1-C02-TH-P050-Y1033-01.webp
% 本文件只在转写完成后作为题干使用；OCR 条目尚不是可靠题干。
\begin{bankitem}{p.50 思考}{第2章 / 2.3 一元二次不等式 / 教材 p.50 / PDF 第57页}{XJ-B1-C02-TH-P050-Y1033}{已对照原页转写题干及公式；待教师复核；未制作答案}
从 2.2 节可以知道，一元二次方程与相应的二次函数有着密切的联系，一元二次方程的根就是对应二次函数的零点。那么，能否利用它们之间的联系，借助二次函数求解一元二次不等式呢？
\end{bankitem}
\end{minipage}\par\medskip
